\FigureLayoutDeclare{img/2010/fujian_science/q12_fig1.png}{11.0cm}{21bp 8bp 35bp 13bp}

\chapter{2010年福建卷（理）}

\section{单选题}

\begin{problem}
计算 \(\sin 43^{\circ} \cos 13^{\circ} - \cos 43^{\circ} \sin 13^{\circ}\) 的结果等于（\quad）

\choices
    {\(\frac{1}{2}\)}
    {\( \frac{\sqrt{3}}{3} \)}
    {\( \frac{\sqrt{2}}{2} \)}
    {\(\frac{\sqrt{3}}{2}\)}
\end{problem}

\begin{answer}
A
\end{answer}

\begin{solution}
由正弦的两角差公式，原式
\[
\sin 43^{\circ}\cos 13^{\circ}-\cos 43^{\circ}\sin 13^{\circ}=\sin 30^{\circ}=\frac{1}{2}
\]
因此选择 A.
\end{solution}

\begin{problem}
以抛物线 \( y^{2} = 4x \) 的焦点为圆心，且过坐标原点的圆的方程为（\quad）

\choices
    {\(x^{2} + y^{2} + 2x = 0\)}
    {\(x^{2} + y^{2} + x = 0\)}
    {\(x^{2} + y^{2} - x = 0\)}
    {\(x^{2} + y^{2} - 2x = 0\)}
\end{problem}

\begin{answer}
D
\end{answer}

\begin{solution}
抛物线 \(y^{2}=4x\) 的标准方程为 \(y^{2}=4px\)，所以 \(p=1\)，焦点为 \(F(1,0)\). 以 \(F\) 为圆心且过原点的圆的半径为 \(1\)，故其方程为
\[
(x-1)^{2}+y^{2}=1
\]
即 \(x^{2}+y^{2}-2x=0\). 因此选择 D.
\end{solution}

\begin{problem}
设等差数列 \(\{a_n\}\) 前 \(n\) 项和为 \(S_n\). 若 \(a_1 = -11\)，\(a_4 + a_6 = -6\)，则当 \(S_n\) 取最小值时，\(n\) 等于（\quad）

\choices
    {6}
    {7}
    {8}
    {9}
\end{problem}

\begin{answer}
A
\end{answer}

\begin{solution}
设等差数列的公差为 \(d\). 由 \(a_{1}=-11\) 及 \(a_{4}+a_{6}=2a_{5}=-6\)，得 \(a_{5}=-3\)，从而
\(d=\frac{a_{5}-a_{1}}{4}=2\)，\(a_{n}=-11+2(n-1)=2n-13\).
于是
\[
S_{n}=\frac{n(a_{1}+a_{n})}{2}=n(n-12)=\left(n-6\right)^{2}-36
\]
当 \(n\) 为正整数时，\(S_{n}\) 在 \(n=6\) 时取得最小值，因此选择 A.
\end{solution}

\begin{problem}
函数 \( f(x) = \begin{cases}
x^2 + 2x - 3, & x \leqslant 0,\\
-2 + \ln x, & x > 0,
\end{cases} \) 的零点个数为（\quad）

\choices
    {0}
    {1}
    {2}
    {3}
\end{problem}

\begin{answer}
C
\end{answer}

\begin{solution}
当 \(x\leqslant0\) 时，令 \(f(x)=0\)，有
\[
x^{2}+2x-3=(x-1)(x+3)=0
\]
其中只有 \(x=-3\) 满足 \(x\leqslant0\). 当 \(x>0\) 时，令 \(f(x)=0\)，得 \(-2+\ln x=0\)，所以 \(x=\e^{2}>0\). 因而函数共有两个零点，选择 C.
\end{solution}

\begin{problem}
阅读如图所示的程序框图，运行相应的程序，输出的 \(i\) 值等于（\quad）

\texfigure[max width=0.40\linewidth]{img_repaint/2010/fujian_science/q5_fig1.tex}

\choices
    {2}
    {3}
    {4}
    {5}
\end{problem}

\begin{answer}
C
\end{answer}

\begin{solution}
程序开始时 \(s=0\)，\(i=1\). 每次循环先计算 \(a=i\times2^{i}\)，再令 \(s=s+a\)，最后令 \(i=i+1\).

第一次循环后，\(a=1\times2^{1}=2\)，\(s=2\)，\(i=2\)，尚有 \(s\leqslant11\). 第二次循环后，\(a=2\times2^{2}=8\)，\(s=10\)，\(i=3\)，仍有 \(s\leqslant11\). 第三次循环后，\(a=3\times2^{3}=24\)，\(s=34>11\)，此时 \(i=4\)，所以输出 \(4\)，选择 C.
\end{solution}

\begin{problem}
如图，若 \(\Omega\) 是长方体 \(ABCD - A_{1}B_{1}C_{1}D_{1}\) 被平面 \(EFGH\) 截去几何体 \(EFGHB_{1}C_{1}\) 后得到的几何体，其中 \(E\) 为线段 \(A_{1}B_{1}\) 上异于 \(B_{1}\) 的点，\(F\) 为线段 \(BB_{1}\) 上异于 \(B_{1}\) 的点，且 \(EH \parallel A_{1}D_{1}\)，则下列结论中不正确的是（\quad）

\bitmapfigure[max width=0.46\linewidth]{img/2010/fujian_science/q6_fig1.png}

\choices
    {\(EH \parallel FG\)}
    {四边形 \(EFGH\) 是矩形}
    {\(\Omega\) 是棱柱}
    {\(\Omega\) 是棱台}
\end{problem}

\begin{answer}
D
\end{answer}

\begin{solution}
由 \(EH\parallel A_{1}D_{1}\parallel BC\)，截面所在平面平行于 \(BC\). 因而该平面在右侧面 \(BCC_{1}B_{1}\) 内的截线 \(FG\) 也平行于 \(BC\)，所以 \(EH\parallel FG\). 另外，\(EF\) 和 \(HG\) 分别是截面与两个相对侧面的交线，故 \(EF\parallel HG\)，从而四边形 \(EFGH\) 是平行四边形. 由于 \(A_{1}D_{1}\perp\) 平面 \(A_{1}ABB_{1}\)，且 \(EH\parallel A_{1}D_{1}\)，而 \(EF\) 在平面 \(A_{1}ABB_{1}\) 内，所以 \(EF\perp EH\)，于是四边形 \(EFGH\) 是矩形.

在平面 \(A_{1}ABB_{1}\) 内，被截去部分的截面为三角形 \(EFB_{1}\)，所以剩余部分的截面为五边形 \(AA_{1}EFB\). 在相对的平面 \(D_{1}DCC_{1}\) 内，剩余部分的截面为与之全等且平行的五边形 \(DD_{1}HGC\)，对应顶点的连线都平行于 \(BC\)，故 \(\Omega\) 是五棱柱. 因此选项 C 正确，选项 D“\(\Omega\) 是棱台”不正确，选择 D.
\end{solution}

\begin{problem}
若点 \(O\) 和点 \(F(-2,0)\) 分别为双曲线 \(\frac{x^2}{a^2} - y^2 = 1\) （\(a > 0\)）的中心和左焦点，点 \(P\) 为双曲线右支上的任意一点，则 \(\overrightarrow{OP} \cdot \overrightarrow{FP}\) 的取值范围为（\quad）

\choices
    {\(\left[3 - 2\sqrt{3}, +\infty\right)\)}
    {\(\left[3 + 2\sqrt{3}, + \infty\right)\)}
    {\(\left[-\frac{7}{4}, +\infty\right)\)}
    {\(\left[\frac{7}{4}, +\infty\right)\)}
\end{problem}

\begin{answer}
B
\end{answer}

\begin{solution}
双曲线的焦距参数满足 \(c=2\)，且 \(c^{2}=a^{2}+1\)，所以 \(a^{2}=3\). 对于右支上的点 \(P(x,y)\)，有 \(x\geqslant\sqrt{3}\)，且
\(\frac{x^{2}}{3}-y^{2}=1\)，\(y^{2}=\frac{x^{2}}{3}-1\).
又 \(\overrightarrow{OP}=(x,y)\)，\(\overrightarrow{FP}=(x+2,y)\)，故
\[
\overrightarrow{OP}\cdot\overrightarrow{FP}=x^{2}+2x+y^{2}=\frac{4}{3}x^{2}+2x-1
\]
令 \(h(x)=\frac{4}{3}x^{2}+2x-1\)，则 \(h'(x)=\frac{8}{3}x+2>0\)（\(x\geqslant\sqrt{3}\)），所以该式在 \(x\geqslant\sqrt{3}\) 上递增，最小值在 \(x=\sqrt{3}\) 时取得，最小值为
\[
\frac{4}{3}\cdot3+2\sqrt{3}-1=3+2\sqrt{3}
\]
所以取值范围为 \(\left[3+2\sqrt{3},+\infty\right)\)，选择 B.
\end{solution}

\begin{problem}
设不等式组 \(\begin{cases}
x \geqslant 1,\\
x - 2y + 3 \geqslant 0,\\
y \geqslant x
\end{cases}\) 所表示的平面区域是 \(\Omega_1\)，平面区域 \(\Omega_2\) 与 \(\Omega_1\) 关于直线 \(3x - 4y - 9 = 0\) 对称. 对于 \(\Omega_1\) 中的任意点 \(A\) 与 \(\Omega_2\) 中的任意点 \(B\)，\(|AB|\) 的最小值等于（\quad）

\choices
    {\(\frac{28}{5}\)}
    {4}
    {\(\frac{12}{5}\)}
    {2}
\end{problem}

\begin{answer}
B
\end{answer}

\begin{solution}
不等式组等价于
\(x\geqslant1\)，\(x\leqslant y\leqslant\frac{x+3}{2}\).
故 \(1\leqslant x\leqslant3\)，区域 \(\Omega_{1}\) 是顶点为 \((1,1)\)，\((1,2)\)，\((3,3)\) 的三角形. 直线 \(3x-4y-9=0\) 在这三个顶点处的左端分别为 \(-10\)，\(-14\)，\(-12\)，所以 \(\Omega_{1}\) 在该直线的一侧，且到直线的最小距离为
\[
\frac{10}{\sqrt{3^{2}+(-4)^{2}}}=2
\]
在点 \((1,1)\) 处取得.

\(\Omega_{2}\) 是 \(\Omega_{1}\) 关于该直线的对称图形. 任取 \(A\in\Omega_{1}\)，\(B\in\Omega_{2}\)，连接 \(A\)，\(B\) 的线段经过该直线时，线段长度不小于 \(A\) 和 \(B\) 到该直线的距离之和，而这两个距离都不小于 \(2\)，故 \(|AB|\geqslant4\). 取 \(A=(1,1)\)，取其关于直线的对称点 \(B\)，则 \(|AB|=4\)，所以最小值为 \(4\)，选择 B.
\end{solution}

\begin{problem}
对于复数 \(a\)，\(b\)，\(c\)，\(d\)，若集合 \(S = \{a, b, c, d\}\) 具有性质“对任意 \(x\)，\(y \in S\)，必有 \(xy \in S\)”，则当 \(\begin{cases}
a = 1,\\
b^2 = 1,\\
c^2 = b
\end{cases}\) 时，\(b + c + d\) 等于（\quad）

\choices
    {1}
    {\(-1\)}
    {0}
    {\(\i\)}
\end{problem}

\begin{answer}
B
\end{answer}

\begin{solution}
因为 \(S\) 是由四个元素组成的集合，\(a\)，\(b\)，\(c\)，\(d\) 互不相同. 由 \(a=1\) 及 \(b^{2}=1\)，且 \(b\neq a\)，得 \(b=-1\). 再由 \(c^{2}=b=-1\)，得 \(c=\i\) 或 \(c=-\i\).

由集合的乘法封闭性，\(bc=-c\in S\). 由于 \(1\)，\(-1\)，\(c\) 已在集合中，第四个元素只能是 \(d=-c\). 因而
\[
b+c+d=-1+c-c=-1
\]
所以选择 B.
\end{solution}

\begin{problem}
对于具有相同定义域 \(D\) 的函数 \(f(x)\) 和 \(g(x)\)，若存在函数 \(h(x) = kx + b\) （\(k\)，\(b\) 为常数），对任给的正数 \(m\)，存在相应的 \(x_0 \in D\)，使得当 \(x \in D\) 且 \(x > x_0\) 时，总有 \(\begin{cases}
0 < f(x) - h(x) < m,\\
0 < h(x) - g(x) < m,
\end{cases}\) 则称直线 \(l: y = kx + b\) 为曲线 \(y = f(x)\) 与 \(y = g(x)\) 的“分渐近线”. 给出定义域均为 \(D = \{x \mid x > 1\}\) 的四组函数如下：
\begin{circlelist}
\item \(f(x)=x^{2}\)，\(g(x) = \sqrt{x}\)；
\item \(f(x) = 10^{-x} + 2\)，\(g(x) = \frac{2x - 3}{x}\)；
\item \(f(x) = \frac{x^2 + 1}{x^2}\)，\(g(x) = \frac{x\ln x + 1}{\ln x}\)；
\item \(f(x) = \frac{2x^2}{x + 1}\)，\(g(x) = 2(x - 1 - \e^{-x})\).
\end{circlelist}
其中，曲线 \( y = f(x) \) 与 \( y = g(x) \) 存在“分渐近线”的是（\quad）
\choices
    {\circled{1}\circled{4}}
    {\circled{2}\circled{3}}
    {\circled{2}\circled{4}}
    {\circled{3}\circled{4}}
\end{problem}

\begin{answer}
C
\end{answer}

\begin{solution}
“分渐近线”的定义等价于存在直线 \(h(x)=kx+b\)，使得当 \(x\to+\infty\) 时，\(f(x)-h(x)\to0^{+}\)，且 \(h(x)-g(x)\to0^{+}\).

对于第\(\circled{1}\)组，无论直线 \(h(x)=kx+b\) 的常数 \(k\)，\(b\) 取何值，都有
\[
x^{2}-h(x)=x^{2}-kx-b\longrightarrow+\infty
\]
不可能满足 \(f(x)-h(x)\to0^{+}\)，所以第\(\circled{1}\)组不满足.

对于第\(\circled{2}\)组，取 \(h(x)=2\)，则
\(f(x)-h(x)=10^{-x}>0\)，\(h(x)-g(x)=2-\frac{2x-3}{x}=\frac{3}{x}>0\)
二者都趋于 \(0\)，所以第\(\circled{2}\)组满足.

对于第\(\circled{3}\)组，\(f(x)=1+\frac{1}{x^{2}}\to1\)，而 \(g(x)=x+\frac{1}{\ln x}\to+\infty\)，不可能满足.

对于第\(\circled{4}\)组，取 \(h(x)=2x-2\)，则
\(f(x)-h(x)=\frac{2}{x+1}>0\)，\(h(x)-g(x)=2\e^{-x}>0\)
二者都趋于 \(0\)，所以第\(\circled{4}\)组满足. 因此满足条件的是第\(\circled{2}\)、\(\circled{4}\)组，选择 C.
\end{solution}

\section{填空题}

\begin{problem}
在等比数列 \(\{a_{n}\}\) 中，若公比 \(q = 4\)，且前3项之和等于21，则该数列的通项公式为 \(a_{n} =\fillinblank{}\).
\end{problem}

\begin{answer}
\(4^{n-1}\)
\end{answer}

\begin{solution}
设首项为 \(a_{1}\). 由前三项之和为 \(21\) 及公比 \(q=4\)，有
\(a_{1}+4a_{1}+16a_{1}=21\)，\(21a_{1}=21\)，\(a_{1}=1\).
所以通项公式为 \(a_{n}=a_{1}q^{n-1}=4^{n-1}\).
\end{solution}

\begin{problem}
若一个底面是正三角形的三棱柱的正视图如图所示，
则其表面积等于\fillinblank{}.
\bitmapfigure[max width=0.40\linewidth]{img/2010/fujian_science/q12_fig1.png}
\end{problem}

\begin{answer}
\(6 + 2\sqrt{3}\)
\end{answer}

\begin{solution}
由正视图可知，该正三棱柱的底面边长为 \(2\)，侧棱长为 \(1\). 因而一个底面的面积为
\[
\frac{1}{2}\times2\times\sqrt{3}=\sqrt{3}
\]
两个底面的总面积为 \(2\sqrt{3}\)，三个侧面的总面积为
\[
3\times2\times1=6
\]
所以三棱柱的表面积为 \(6+2\sqrt{3}\).
\end{solution}

\begin{problem}
某次知识竞赛规则如下：在主办方预设的5个问题中，选手若能连续正确回答出两个问题，即停止答题，晋级下一轮. 假设某选手正确回答每个问题的概率都是0.8，且每个问题的回答结果相互独立，则该选手恰好回答了4个问题就晋级下一轮的概率等于\fillinblank{}.
\end{problem}

\begin{answer}
0.128
\end{answer}

\begin{solution}
记回答正确为 \(C\)，回答错误为 \(W\). 要恰好回答到第4个问题才因连续正确两个问题而晋级，第3、4个问题必须为 \(C\)，且第2个问题必须为 \(W\)，否则第2、3个问题已经连续正确. 第1个问题可为 \(C\) 或 \(W\)，因此所求概率为
\[
(0.8+0.2)\times0.2\times0.8\times0.8=0.128
\]
\end{solution}

\begin{problem}
已知函数 \( f(x) = 3\sin \left( \omega x - \frac{\pi}{6} \right) \)（\(\omega > 0\)）和 \( g(x) = 2\cos (2x + \varphi) + 1 \) 的图象的对称轴完全相同. 若 \( x \in \left[0, \frac{\pi}{2}\right] \)，则 \( f(x) \) 的取值范围是\fillinblank{}.
\end{problem}

\begin{answer}
\(\left[-\frac{3}{2},3\right]\)
\end{answer}

\begin{solution}
函数 \(f\) 的对称轴为
\(x=\frac{\frac{\pi}{2}+k\pi+\frac{\pi}{6}}{\omega}=\frac{\frac{2\pi}{3}+k\pi}{\omega}\)，\(k\in\Z\).
函数 \(g\) 的对称轴为 \(2x+\varphi=k\pi\)，相邻对称轴的距离为 \(\frac{\pi}{2}\). 两组对称轴完全相同，所以 \(\frac{\pi}{\omega}=\frac{\pi}{2}\)，即 \(\omega=2\). 因而在 \(0\leqslant x\leqslant\frac{\pi}{2}\) 时，
\[
f(x)=3\sin\left(2x-\frac{\pi}{6}\right)
\]
其自变量范围为 \(\left[-\frac{\pi}{6},\frac{5\pi}{6}\right]\). 该范围包含 \(\frac{\pi}{2}\)，所以最大值为 \(3\)；左端点给出最小值 \(3\sin\left(-\frac{\pi}{6}\right)=-\frac{3}{2}\). 因此取值范围为 \(\left[-\frac{3}{2},3\right]\).
\end{solution}

\begin{problem}
已知定义域为 \((0, +\infty)\) 的函数 \(f(x)\) 满足：
\begin{enumerate}
\item 对任意 \(x \in (0, +\infty)\)，恒有 \(f(2x) = 2f(x)\) 成立；
\item 当 \(x \in (1, 2]\) 时，\(f(x) = 2 - x\).
\end{enumerate}
给出结论如下：
\begin{circlelist}
\item 对任意 \( m \in \Z \)，有 \( f(2^{m}) = 0 \)；
\item 函数 \( f(x) \) 的值域为 \( [0, +\infty) \)；
\item 存在 \( n \in \Z \)，使得 \( f(2^{n+1}) = 9 \)；
\item “函数 \( f(x) \) 在区间 \( (a, b) \) 上单调递减”的充要条件是“存在 \( k \in \Z \)，使得 \( (a, b) \subseteq (2^{k}, 2^{k+1}) \)”.
\end{circlelist}
其中所有正确结论的序号是\fillinblank{}.
\end{problem}

\begin{answer}
\(\circled{1}\circled{2}\circled{4}\)
\end{answer}

\begin{solution}
由 \(f(2)=0\) 及 \(f(2)=2f(1)\)，得 \(f(1)=0\). 对任意整数 \(m\)，反复使用 \(f(2x)=2f(x)\) 及其等价式 \(f(x)=2f(x/2)\)，可得 \(f(2^{m})=0\)，所以结论\(\circled{1}\)正确.

对任意整数 \(k\)，当 \(2^{k}<x\leqslant2^{k+1}\) 时，令 \(u=2^{-k}x\)，则 \(1<u\leqslant2\)，从而
\[
f(x)=2^{k}f(u)=2^{k}(2-u)=2^{k+1}-x
\]
在每个区间 \(\left(2^{k},2^{k+1}\right]\) 上，函数值属于 \([0,2^{k})\)，而 \(k\) 可任意增大，所以值域为 \([0,+\infty)\)，结论\(\circled{2}\)正确.

由结论\(\circled{1}\)，所有 \(f(2^{n+1})\) 都等于 \(0\)，不可能等于 \(9\)，所以结论\(\circled{3}\)错误.

在每个开区间 \(\left(2^{k},2^{k+1}\right)\) 内，\(f(x)=2^{k+1}-x\)，严格单调递减. 若 \((a,b)\) 跨过某个 \(2^{k}\)，则函数在该点右侧由 \(0\) 跳到接近 \(2^{k}\) 的正值，不可能在整个 \((a,b)\) 上单调递减. 因此结论\(\circled{4}\)正确. 所有正确结论的序号为 \(\circled{1}\)\(\circled{2}\)\(\circled{4}\).
\end{solution}

\section{解答题}

\begin{problem}
设 \(S\) 是不等式 \(x^{2} - x - 6 \leqslant 0\) 的解集，整数 \(m\)，\(n \in S\).
\begin{enumerate}
\item 记使得“\(m + n = 0\) 成立的有序数组 \((m, n)\)”为事件 \(A\)，试列举 \(A\) 包含的基本事件；
\item 设 \(\xi = m^2\)，求 \(\xi\) 的分布列及其数学期望 \(E(\xi)\).
\end{enumerate}
\end{problem}

\begin{solution}
\begin{enumerate}
\item 解不等式得
\[
x^{2}-x-6=(x-3)(x+2)\leqslant0
\]
所以 \(S=[-2,3]\). 因为 \(m\)，\(n\) 是 \(S\) 中的整数，故 \(m\)，\(n\) 各有 \(6\) 种等可能取值，有序数组共有 \(36\) 个.

事件 \(A\) 要求 \(n=-m\). 在 \(S\) 中逐一取 \(m=-2\)，\(-1\)，\(0\)，\(1\)，\(2\)，\(3\)，得到
\[
A=\{(-2,2),(-1,1),(0,0),(1,-1),(2,-2)\}
\]

\item 由于题目要求计算概率和分布列，按随机试验的标准解释，\(m\)，\(n\) 在 \(S\) 中的整数里独立等可能取值，故 \(m\) 在 \(\{-2,-1,0,1,2,3\}\) 中等可能取值，
\(P(\xi=0)=\frac16\)，\(P(\xi=1)=\frac26=\frac13\)
\(P(\xi=4)=\frac26=\frac13\)，\(P(\xi=9)=\frac16\).
因此
\[
E(\xi)=0\cdot\frac16+1\cdot\frac13+4\cdot\frac13+9\cdot\frac16=\frac{19}{6}
\]
\end{enumerate}
\end{solution}

\begin{problem}
已知中心在坐标原点 \(O\) 的椭圆 \(C\) 经过点 \(A(2,3)\)，且点 \(F(2,0)\) 为其右焦点.
\begin{enumerate}
\item 求椭圆 \(C\) 的方程；
\item 是否存在平行于 \(OA\) 的直线 \(l\)，使得直线 \(l\) 与椭圆 \(C\) 有公共点，且直线 \(OA\) 与 \(l\) 的距离等于 \(4\)？若存在，求出直线 \(l\) 的方程；若不存在，请说明理由.
\end{enumerate}
\end{problem}

\begin{solution}
\begin{enumerate}
\item 设椭圆的标准方程为
\(\frac{x^{2}}{a^{2}}+\frac{y^{2}}{b^{2}}=1\)，\(a>b>0\).
右焦点为 \(F(2,0)\)，所以 \(c=2\)，\(a^{2}-b^{2}=c^{2}=4\). 由点 \(A(2,3)\) 在椭圆上，得
\[
\frac{4}{a^{2}}+\frac{9}{b^{2}}=1
\]
令 \(a^{2}=u\)，则 \(b^{2}=u-4\)，从而
\(\frac4u+\frac9{u-4}=1\)，\(\Longrightarrow\)，\(u^{2}-17u+16=0\).
解得 \(u=1\) 或 \(u=16\). 因 \(u=a^{2}>c^{2}=4\)，故 \(a^{2}=16\)，\(b^{2}=12\)，椭圆方程为
\[
\frac{x^{2}}{16}+\frac{y^{2}}{12}=1
\]

\item 直线 \(OA\) 的方程为 \(3x-2y=0\). 设与它平行的直线为
\[
l:3x-2y+\lambda=0
\]
由两平行直线的距离为 \(4\)，得
\(\frac{|\lambda|}{\sqrt{13}}=4\)，\(\lambda^{2}=208\).
在直线 \(l\) 上，\(y=(3x+\lambda)/2\). 代入椭圆方程，得
\[
3x^{2}+(3x+\lambda)^{2}=48
\]
即
\[
12x^{2}+6\lambda x+\lambda^{2}-48=0
\]
其判别式为
\[
\Delta=(6\lambda)^{2}-4\times12(\lambda^{2}-48)=2304-12\lambda^{2}=-192<0
\]
所以直线 \(l\) 与椭圆没有公共点，不存在这样的直线.
\end{enumerate}
\end{solution}

\begin{problem}
如图，圆柱 \(O O_{1}\) 内有一个三棱柱 \(A B C - A_{1}B_{1}C_{1}\)，三棱柱的底面为圆柱底面的内接三角形，且 \(A B\) 是圆 \(O\) 的直径.

\begin{enumerate}
\item 证明：平面 \(A_{1}ACC_{1}\) \(\perp\) 平面 \(B_{1}BCC_{1}\)；
\item 设 \(AB = AA_{1}\)，在圆柱 \(OO_{1}\) 内随机选取一点. 记该点取自于三棱柱 \(ABC - A_{1}B_{1}C_{1}\) 内的概率为 \(p\).
\begin{enumerate}
\item 当点 \(C\) 在圆周上运动时，求 \(p\) 的最大值；
\item 记平面 \(A_{1}ACC_{1}\) 与平面 \(B_{1}OC\) 所成的角为 \(\theta\) （\(0^{\circ} < \theta \leqslant 90^{\circ}\)），当 \(p\) 取最大值时，求 \(\cos \theta\) 的值.
\end{enumerate}
\end{enumerate}

\bitmapfigure[max width=0.46\linewidth]{img/2010/fujian_science/q18_fig1.png}
\end{problem}

\begin{solution}
\begin{enumerate}
\item 因为 \(AB\) 是圆 \(O\) 的直径，根据圆周角定理，
\(\angle ACB=90^{\circ}\)，\(AC\perp BC\).
由图中三棱柱的侧棱垂直于底面，得 \(CC_{1}\perp\) 平面 \(ABC\). 平面 \(A_{1}ACC_{1}\) 与平面 \(B_{1}BCC_{1}\) 的交线为 \(CC_{1}\)，它们分别在底面 \(ABC\) 上的截线为 \(AC\) 和 \(BC\)，且 \(AC\perp CC_{1}\)，\(BC\perp CC_{1}\). 因此两平面所成的二面角等于其在底面上的截线 \(AC\) 与 \(BC\) 的夹角. 由于 \(AC\perp BC\)，故两平面所成的二面角为直角，即
\[
\text{平面 }A_{1}ACC_{1}\perp\text{平面 }B_{1}BCC_{1}
\]

\item 设圆柱底面半径为 \(r\). 则 \(AB=2r\)，又 \(AB=AA_{1}\)，所以圆柱高为 \(2r\). 三棱柱与圆柱的高相同，故
\[
p=\frac{\frac12 AC\cdot BC\cdot2r}{\pi r^{2}\cdot2r}=\frac{AC\cdot BC}{2\pi r^{2}}
\]
由 \(\angle ACB=90^{\circ}\)，有
\[
AC^{2}+BC^{2}=AB^{2}=4r^{2}
\]
因此
\[
AC\cdot BC\leqslant\frac{AC^{2}+BC^{2}}{2}=2r^{2}
\]
当且仅当 \(AC=BC=\sqrt{2}r\) 时取等号. 所以
\[
p_{\max}=\frac1\pi
\]

当 \(p\) 取最大值时，取空间直角坐标系，使
\(O=(0,0,0)\)，\(A=(-r,0,0)\)，\(B=(r,0,0)\)，\(C=(0,r,0)\)
\(A_{1}=(-r,0,2r)\)，\(B_{1}=(r,0,2r)\).
平面 \(A_{1}ACC_{1}\) 的一个法向量可取 \(\bs{n}_{1}=(1,-1,0)\)，平面 \(B_{1}OC\) 的一个法向量可取 \(\bs{n}_{2}=(2,0,-1)\). 因而
\[
\cos\theta=\frac{|\bs{n}_{1}\cdot\bs{n}_{2}|}{|\bs{n}_{1}|\,|\bs{n}_{2}|}
=\frac{2}{\sqrt{2}\sqrt{5}}=\frac{\sqrt{10}}{5}
\]
\end{enumerate}
\end{solution}

\begin{problem}
某港口 \(O\) 要将一件重要物品用小艇送到一艘正在航行的轮船上. 在小艇出发时，轮船位于港口 \(O\) 北偏西 \(30^{\circ}\) 且与该港口相距 \(20\text{ 海里}\) 的 \(A\) 处，并以 \(30\text{ 海里/小时}\) 的航行速度沿正东方向匀速行驶. 假设该小艇沿直线方向以 \(v\text{ 海里/小时}\) 的航行速度匀速行驶，经过 \(t\text{ 小时}\) 与轮船相遇.
\begin{enumerate}
\item 若希望相遇时小艇的航行距离最小，则小艇航行速度的大小应为多少？
\item 假设小艇的最高航行速度只能达到 \(30\text{ 海里/小时}\)，试设计航行方案（即确定航行方向与航行速度的大小），使得小艇能以最短时间与轮船相遇，并说明理由.
\end{enumerate}
\end{problem}

\begin{solution}
\begin{enumerate}
\item 取港口 \(O\) 为原点，正东方向为 \(x\) 轴，正北方向为 \(y\) 轴. 轮船初始位置为
\[
A=(-10,10\sqrt{3})
\]
经过 \(t\) 小时后的位置为
\[
R(t)=(30t-10,10\sqrt{3})
\]
小艇航行距离的平方为
\[
d^{2}(t)=(30t-10)^{2}+300=900\left(t-\frac13\right)^{2}+300
\]
所以当 \(t=\frac13\) 时距离最小，最小距离为 \(10\sqrt3\). 此时
\[
v=\frac{d(t)}{t}=\frac{10\sqrt3}{1/3}=30\sqrt3\text{ 海里/小时}
\]

\item 若小艇速度不超过 \(30\text{ 海里/小时}\)，相遇必须满足 \(d(t)\leqslant30t\). 因而
\((30t-10)^{2}+300\leqslant900t^{2}\)，\(\Longrightarrow\)，\(t\geqslant\frac23\).
所以最短相遇时间为 \(t=\frac23\). 此时轮船位置为
\[
R\left(\frac23\right)=(10,10\sqrt3)
\]
小艇从 \(O\) 到该点的距离为 \(20\)，所需速度为
\[
\frac{20}{2/3}=30\text{ 海里/小时}
\]
方向向量为 \((10,10\sqrt3)\)，即与正东方向成 \(60^{\circ}\) 角，也就是北偏东 \(30^{\circ}\). 该方案达到速度上限且时间最短.
\end{enumerate}
\end{solution}

\begin{problem}
\begin{enumerate}
\item 已知函数 \( f(x) = x^3 - x \)，其图象记为曲线 \(C\).
    \begin{enumerate}
    \item 求函数 \( f(x) \) 的单调区间；
    \item 证明：若对于任意非零实数 \(x_{1}\)，曲线 \(C\) 与其在点 \(P_{1}(x_{1}, f(x_{1}))\) 处的切线交于另一点 \(P_{2}(x_{2}, f(x_{2}))\)，曲线 \(C\) 与其在点 \(P_{2}(x_{2}, f(x_{2}))\) 处的切线交于另一点 \(P_{3}(x_{3}, f(x_{3}))\)，线段 \(P_{1}P_{2}\)，\(P_{2}P_{3}\) 与曲线 \(C\) 所围成封闭图形的面积分别为 \(S_{1}\)，\(S_{2}\)，则 \(\frac{S_{1}}{S_{2}}\) 为定值；
    \end{enumerate}
\item 对于一般的三次函数 \( g(x) = ax^3 + bx^2 + cx + d \) （\(a \neq 0\)），请给出类似于（1）（ii） 的正确命题，并予以证明.
\end{enumerate}
\end{problem}

\begin{solution}
\begin{enumerate}
\item
\begin{enumerate}
\item \(f'(x)=3x^{2}-1\). 当 \(x<-\frac{\sqrt3}{3}\) 或 \(x>\frac{\sqrt3}{3}\) 时，\(f'(x)>0\)；当 \(-\frac{\sqrt3}{3}<x<\frac{\sqrt3}{3}\) 时，\(f'(x)<0\). 因此函数在
\[
\left(-\infty,-\frac{\sqrt3}{3}\right)\quad\text{和}\quad\left(\frac{\sqrt3}{3},+\infty\right)
\]
上单调递增，在
\[
\left(-\frac{\sqrt3}{3},\frac{\sqrt3}{3}\right)
\]
上单调递减.

\item 设曲线在 \(P_{1}(x_{1},f(x_{1}))\) 处的切线为 \(L_{1}\). 因为 \(f'(x_{1})=3x_{1}^{2}-1\)，所以
\[
L_{1}:y=(3x_{1}^{2}-1)x-2x_{1}^{3}
\]
曲线与该切线的交点横坐标满足
\[
x^{3}-x=(3x_{1}^{2}-1)x-2x_{1}^{3}
\]
即
\[
(x-x_{1})^{2}(x+2x_{1})=0
\]
除去切点 \(x=x_{1}\)，另一交点的横坐标为 \(x_{2}=-2x_{1}\). 同理，\(x_{3}=-2x_{2}=4x_{1}\).

切线与曲线之差为
\[
f(x)-L_{1}(x)=(x-x_{1})^{2}(x+2x_{1})
\]
令 \(x=x_{1}+(x_{2}-x_{1})u\)，可得两交点之间所围面积为
\[
S_{1}=|x_{2}-x_{1}|^{4}\int_{0}^{1}u^{2}(1-u)\,\mathrm du
=\frac{|x_{2}-x_{1}|^{4}}{12}
\]
同理
\[
S_{2}=\frac{|x_{3}-x_{2}|^{4}}{12}
\]
由于 \(x_{2}=-2x_{1}\)，\(x_{3}=-2x_{2}\)，有 \(|x_{3}-x_{2}|=2|x_{2}-x_{1}|\). 因此
\[
\frac{S_{1}}{S_{2}}=\frac{1}{2^{4}}=\frac{1}{16}
\]
为定值.
\end{enumerate}
\item 设一般三次函数为 \(g(x)=ax^{3}+bx^{2}+cx+d\)，其中 \(a\neq0\). 正确命题为：对于任意满足 \(3ax_{1}+b\neq0\) 的实数 \(x_{1}\)，若曲线在 \(P_{1}\) 处的切线与曲线交于另一点 \(P_{2}\)，再在 \(P_{2}\) 处作切线并与曲线交于另一点 \(P_{3}\)，则由线段 \(P_{1}P_{2}\)、\(P_{2}P_{3}\) 及曲线围成的面积 \(S_{1}\)、\(S_{2}\) 满足
\[
\frac{S_{1}}{S_{2}}=\frac{1}{16}
\]

事实上，设切点横坐标为 \(u\)，其切线为 \(L_{u}\). 直接展开得
\[
g(x)-L_{u}(x)=(x-u)^{2}\left(a(x+2u)+b\right)
\]
所以另一交点的横坐标为
\[
v=-2u-\frac{b}{a}
\]
条件 \(3au+b\neq0\) 正好保证 \(v\neq u\). 对下一点有
\(w=-2v-\frac{b}{a}\)，\(|w-v|=2|v-u|\).
而在横坐标为 \(u\)，\(v\) 的两个交点之间，面积为
\[
|a|\int_{\min\{u,v\}}^{\max\{u,v\}}|(x-u)^{2}(x-v)|\,\mathrm dx
=\frac{|a|}{12}|v-u|^{4}
\]
因此
\[
\frac{S_{1}}{S_{2}}=\frac{|v-u|^{4}}{|w-v|^{4}}=\frac{1}{16}
\]
参数 \(c\)，\(d\) 在切线差 \(g(x)-L_u(x)\) 中消去，因此不影响上述比值.
\end{enumerate}
\end{solution}

\section{选考题：解答题}

\begin{problem}
已知矩阵 \(M = \begin{pmatrix}1 & a\\ b & 1 \end{pmatrix}\)，\(N = \begin{pmatrix}c & 2\\ 0 & d \end{pmatrix}\)，且 \(MN = \begin{pmatrix}2 & 0\\ -2 & 0 \end{pmatrix}\).
\begin{enumerate}
\item 求实数 \(a\)，\(b\)，\(c\)，\(d\) 的值；
\item 求直线 \(y = 3x\) 在矩阵 \(M\) 所对应的线性变换下的像的方程.
\end{enumerate}
\end{problem}

\begin{solution}
\begin{enumerate}
\item 由矩阵乘法
\[
MN=\begin{pmatrix}c&2+ad\\bc&2b+d\end{pmatrix}=\begin{pmatrix}2&0\\-2&0\end{pmatrix}
\]
得到
\(c=2\)，\(ad=-2\)，\(bc=-2\)，\(2b+d=0\).
由 \(c=2\) 和 \(bc=-2\) 得 \(b=-1\)，再由 \(2b+d=0\) 得 \(d=2\)，最后由 \(ad=-2\) 得 \(a=-1\). 所以
\(a=-1\)，\(b=-1\)，\(c=2\)，\(d=2\).

\item 矩阵 \(M\) 对点 \((x,y)\) 的变换为
\[
\begin{pmatrix}x'\\y'\end{pmatrix}
=\begin{pmatrix}1&-1\\-1&1\end{pmatrix}
\begin{pmatrix}x\\y\end{pmatrix},
\qquad x'=x-y,\quad y'=-x+y
\]
原直线上点可写为 \((t,3t)\). 其像为
\[
(x',y')=(-2t,2t)
\]
所以像上的点满足 \(y'=-x'\). 又 \(t\) 可取任意实数，故直线的像为 \(y=-x\).
\end{enumerate}
\end{solution}

\begin{problem}
在直角坐标系 \(xOy\) 中，直线 \(l\) 的参数方程为 \(\begin{cases}
x = 3 - \frac{\sqrt{2}}{2} t,\\
y = \sqrt{5} - \frac{\sqrt{2}}{2} t.
\end{cases}\)（\(t\) 为参数）. 在极坐标系（与直角坐标系 \(xOy\) 取相同的长度单位，且以原点 \(O\) 为极点，以 \(x\) 轴正半轴为极轴）中，圆 \(C\) 的方程为 \(\rho = 2\sqrt{5}\sin \theta\).
\begin{enumerate}
\item 求圆 \(C\) 的直角坐标方程；
\item 设圆 \(C\) 与直线 \(l\) 交于点 \(A\)，\(B\)，若点 \(P\) 的坐标为 \((3, \sqrt{5})\)，求 \(|PA| + |PB|\).
\end{enumerate}
\end{problem}

\begin{solution}
\begin{enumerate}
\item 由极坐标与直角坐标的关系 \(x=\rho\cos\theta\)，\(y=\rho\sin\theta\)，有
\(\rho^{2}=x^{2}+y^{2}\)，\(\rho\sin\theta=y\).
将 \(\rho=2\sqrt5\sin\theta\) 两边乘以 \(\rho\)，得
\[
x^{2}+y^{2}=2\sqrt5y
\]
所以圆 \(C\) 的直角坐标方程为
\[
x^{2}+y^{2}-2\sqrt5y=0
\]

\item 直线上的点可参数化为
\(x=3-\frac{\sqrt2}{2}t\)，\(y=\sqrt5-\frac{\sqrt2}{2}t\).
将其代入圆的方程，得
\[
t^{2}-3\sqrt2t+4=0
\]
所以两个交点对应的参数为
\(t_{1}=\sqrt2\)，\(t_{2}=2\sqrt2\).
参数方程中的方向向量长度为 \(1\)，且 \(t=0\) 对应点 \(P(3,\sqrt5)\)，故
\(|PA|=|t_{1}|=\sqrt2\)，\(|PB|=|t_{2}|=2\sqrt2\).
因此 \(|PA|+|PB|=3\sqrt2\).
\end{enumerate}
\end{solution}

\begin{problem}
已知函数 \(f(x) = |x - a|\).
\begin{enumerate}
\item 若不等式 \(f(x) \leqslant 3\) 的解集为 \(\{x \mid -1 \leqslant x \leqslant 5\}\)，求实数 \(a\) 的值；
\item 在（1）的条件下，若 \(f(x) + f(x + 5) \geqslant m\) 对一切实数 \(x\) 恒成立，求实数 \(m\) 的取值范围.
\end{enumerate}
\end{problem}

\begin{solution}
\begin{enumerate}
\item 不等式 \(|x-a|\leqslant3\) 的解集为
\[
a-3\leqslant x\leqslant a+3
\]
与题给解集 \([-1,5]\) 比较端点，得 \(a-3=-1\)，所以 \(a=2\).

\item 代入 \(a=2\)，有
\[
f(x)+f(x+5)=|x-2|+|x+3|
\]
由三角不等式，
\[
|x-2|+|x+3|\geqslant|(x-2)-(x+3)|=5
\]
当 \(-3\leqslant x\leqslant2\) 时等号成立，所以左端的最小值为 \(5\). 因而不等式对一切实数 \(x\) 恒成立，当且仅当 \(m\leqslant5\).
\end{enumerate}
\end{solution}
